%%=====================================================================
%%  附录
%%  【规范】附录：（另起一页，小 4 号黑体，顶格）；附录中的公式编号为（A1）。
%%=====================================================================

\begin{theappendix}

本附录给出系统实现过程中的部分核心代码，供评阅与复现参考。

\subsection*{A1\quad 首页视图层结构}

\begin{lstlisting}[style=xml, caption={首页视图层 index.wxml}]
<view class="container" style='align-items: center;'>
  <view class='top_notice' wx:if="{{notAuthorized}}">
    <text>小程序需要您的授权才能正常使用</text>
    <button bindtap='onOpenSetting'>去授权</button>
  </view>
  <view class='index_card_img' bindtap='toPunchcardPage'
        style="background-image:url(https://xxx.com/card.png);">
  </view>
  <view class='index_text_bg'>
    <text class='index_title'>幸运占卜</text>
    <text class='index_subtitle'>每日需刷满20题可进行一次占卜</text>
  </view>
  <view class='countDown'>
    <image src='/images/punchcard/{{day.left}}.png' class='l'></image>
    <image src='/images/punchcard/{{day.right}}.png' class='r'></image>
  </view>
</view>
\end{lstlisting}

\subsection*{A2\quad 抽卡动画处理逻辑}

\begin{lstlisting}[style=js, caption={洗牌动画 onShuffleEvent}]
/* 洗牌 */
onShuffleEvent: function (cb) {
  shuffleCount++;
  this.innerAudioContext1.play();          // 洗牌音效
  this.shuffleAnimation2.rotateY(0).step();
  wsAPI.taskSequence()
    .then(() => {
      this.shuffleAnimation1.left(-70).step();
      this.shuffleAnimation2.left(70).step();
    })
    .then(() => {
      var that = this;
      clearTimeout(timer);
      var timer = setTimeout(function () { cb && cb(that); }, 300);
    });
}
\end{lstlisting}

\subsection*{A3\quad 每日可占卜次数判定}

用户当日可占卜次数由当日累计答题数与占卜基数共同决定，其判定公式为

\begin{equation}
  C = \left\lfloor \frac{Q}{20} \right\rfloor - U
  \label{equ:appendix-count}
\end{equation}

式中：$C$ 为当日剩余可占卜次数；$Q$ 为当日累计答题数；$U$ 为当日已使用次数；$\lfloor \cdot \rfloor$ 表示向下取整。

\end{theappendix}
